
\section*{Problématique : Comment enseigner une langue « morte » comme une langue vivante pour la sauver ?}

C'est de cette tension que naît la question centrale de cette thèse : \textbf{Comment enseigner le grec ancien comme on enseigne une langue vivante, non pas malgré son statut de langue historique, mais en l'exploitant comme une ressource ?}

Il ne s'agit pas de renier la rigueur historique. Mais il s'agit de reconnaître une évidence oubliée : le grec ancien n'est « mort » que pour nous, pas en lui-même. Il survit dans le grec moderne (continuité lexicale à 70-80\%). Il survit dans le vocabulaire scientifique international (80-90\% des termes médicaux, biologiques, technologiques sont d'origine gréco-latine). Il survit dans la culture vivante de la Grèce contemporaine.

Une refondation cognitive de l'enseignement du grec ancien doit partir d'une double stratégie :
\begin{enumerate}
\item \textbf{Priorité lexicale} : Abandon de la grammaire comme socle pour la remplacer par un apprentissage massif du noyau lexical fondamental (1100 lemmes = 80\% de couverture textuelle).
\item \textbf{Réactivation de la continuité} : Utilisation du grec moderne et du vocabulaire scientifique comme points d'ancrage motivationnels et cognitifs.
\end{enumerate}

Cette approche n'est pas nouvelle en didactique des langues (elle s'appelle \textit{Second Language Acquisition} ou SLA). Mais elle n'a jamais été systématiquement appliquée au grec ancien en France. Des institutions comme l'Institut Polis en Israël et Elliniki Agogi en Grèce le démontrent : on peut transformer le grec ancien en une langue vivante d'apprentissage rapide et de communication réelle.

\vspace{0.5cm}

\noindent\rule{\textwidth}{0.5pt}

Cette thèse se structure en trois mouvements :

\textbf{I. Partie Historique.} Elle retraçe l'évolution longue durée de l'enseignement du grec ancien de 1870 à 2024, mettant au jour comment un régime pédagogique fondé sur l'analyse philologique, l'élitisme et la croyance en le transfert cognitif s'est cristallisé, maintenu et finalement sclérosé. Elle montre comment les réformes institutionnelles successives (1902, 1975, 2016) n'ont jamais remis en cause le paradigme fondamental hérité de 1880.

\textbf{II. Partie Critique.} Elle déploie les résultats des sciences cognitives, de la linguistique acquisitionnelle et des neurosciences pour invalider les présupposés pédagogiques de la Partie I. Elle démontre, preuves à l'appui, que le régime philologique-grammatical produit l'effet inverse de celui qu'il prétend produire.

\textbf{III. Partie Constructive.} Elle propose un protocole pédagogique alternatif fondé sur la priorité lexicale, l'approche communicative et l'activation de la continuité greco-moderne. Elle explore le rôle de la prononciation, du grec contemporain et du vocabulaire scientifique comme leviers de motivation et d'efficacité acquisitionnelle.
